% Part: lambda-calculus
% Chapter: syntax
% Section: beta

\documentclass[../../../include/open-logic-section]{subfiles}

\begin{document}

\olfileid{lam}{syn}{bet}
\olsection{Reduksi-$\beta$}

Nalika ndeleng
$(\lambd[m][(\lambd[y][y]) m])$,
lumrah yen kita ngira ana
sesambungan karo $\lambd[m][m]$:
suku kapindho iku asil
``nyederhanakake'' suku kapisan.
Konsep \emph{reduksi-$\beta$}
ngandharake intuisi iki kanthi formal.

\begin{defn}[Kontraksi-$\beta$, $\bredone$] \ollabel{defn:betacontr}
  \emph{Kontraksi-$\beta$} ($\bredone$)
  yaiku relasi kompatibel paling cilik
  ing suku kang nyukupi syarat iki:
  \[
    (\lambd[x][N])Q \bredone \Subst{N}{Q}{x} 
  \]
  $P$ diarani \emph{dikontraksi-$\beta$}
  dadi $Q$ yen $P \bredone Q$.
  Suku kang awujud
  $(\lambd[x][N])Q$ diarani
  \emph{redex}.
\end{defn}

\begin{prob} \ollabel{prob:def}
  Tulisen definisi induktif
  kang ekuivalen kanggo
  kontraksi-$\beta$, kaya
  owah-owahan variabel kaiket
  ing \olref[lam][syn][alp]{defn:aconvone}.
\end{prob}
  
\begin{defn}[Reduksi-$\beta$, $\bred$] \ollabel{defn:betared}
  \emph{Reduksi-$\beta$} ($\bred$)
  yaiku relasi refleksif lan transitif
  paling cilik ing suku kang
  ngemot $\bredone$.
  $P$ diarani direduksi-$\beta$
  dadi $Q$ yen $P \bred Q$.
\end{defn}

Yen konteks wis cetha,
kita nulis $\redone$ tinimbang
$\bredone$, lan $\red$
tinimbang $\bred$.

Miturut intuisi, $M \bred N$
yen lan mung yen $M$ bisa
diowahi dadi $N$ kanthi nol
utawa sawetara langkah
kontraksi-$\beta$.

\begin{defn}[Normal-$\beta$]
Suku kang ora bisa dikontraksi-$\beta$
maneh diarani \emph{normal-$\beta$}.
\end{defn}

Yen $M \bred N$ lan $N$
normal-$\beta$, mula $N$
diarani \emph{wujud normal}
saka~$M$. Apa wujud normal
sawijining suku iku tunggal?
Ya, kaya kang bakal kita
deleng sabanjure.

Gatekna sawetara conto.
\begin{enumerate}
\item Kita nduweni
\begin{align*}
(\lambd[x][xxy]) \lambd[z][z] & \redone (\lambd[z][z])(\lambd[z][z]) y \\
& \redone (\lambd[z][z]) y \\
& \redone y
\end{align*}
\item ``Nyederhanakake'' suku
  malah bisa ndadekake suku
  luwih rumit:
\begin{align*}
(\lambd[x][xxy])(\lambd[x][xxy]) & \redone (\lambd[x][xxy])(\lambd[x][xxy])y \\
& \redone (\lambd[x][xxy])(\lambd[x][xxy])yy \\
& \redone \dots
\end{align*}
\item Reduksi uga bisa
  ora ngowahi sawijining suku:
\[
(\lambd[x][xx])(\lambd[x][xx]) \redone (\lambd[x][xx])(\lambd[x][xx])
\]
\item Sawetara suku uga bisa
  direduksi kanthi luwih saka siji
  cara; umpamane,
\[
(\lambd[x][(\lambd[y][yx]) z]) v \redone (\lambd[y][yv]) z
\]
kanthi ngontraksi aplikasi
paling njaba; lan
\[
(\lambd[x][(\lambd[y][yx]) z]) v \redone (\lambd[x][zx]) v
\]
kanthi ngontraksi aplikasi
paling njero. Nanging gatekna
yen ing kasus iki loro-lorone
suku terus direduksi dadi suku
kang padha, yaiku $zv$.
\end{enumerate}

Asil pungkasan ing conto mau
dudu kabeneran, nanging
nggambarake sipat kalkulus lambda
kang jero lan wigati,
yaiku sipat Church--Rosser.

\begin{digress}
  Umume ana luwih saka siji
  cara kanggo ngreduksi-$\beta$
  sawijining suku. Mula ana
  maneka strategi reduksi;
  kang paling lumrah yaiku
  \emph{strategi alami}.
  Strategi alami tansah ngontraksi
  redex \emph{paling kiwa},
  kanthi posisi redex ditetepake
  minangka titik wiwitane
  ing suku. Strategi alami nduweni
  sipat migunani: sawijining suku
  bisa direduksi dadi wujud normal
  kanthi sawijining strategi
  yen lan mung yen bisa direduksi
  dadi wujud normal kanthi
  strategi alami. Sabanjure
  kita nggunakake strategi alami
  kajaba yen ditemtokake liyane.
\end{digress}

\begin{defn}[Ekuivalensi-$\beta$, $\equal$]
  \emph{Ekuivalensi-$\beta$} ($\equal$)
  yaiku relasi kang ditetepake
  kanthi induktif kaya mangkene:
  \begin{enumerate}
  \item $M \equal M$.
  \item Yen $M \equal N$, mula $N \equal M$.
  \item Yen $M \equal N$, $N \equal O$, mula $M \equal O$.
  \item Yen $M \equal N$, mula $PM \equal PN$.
  \item Yen $M \equal N$, mula $MQ \equal NQ$.
  \item Yen $M \equal N$, mula $\lambd[x][M] \equal \lambd[x][N]$.
  \item $(\lambd[x][N])Q \equal \Subst{N}{Q}{x}$.
  \end{enumerate}
\end{defn}

Telung aturan kapisan ndadekake
relasi iki relasi ekuivalensi;
telung aturan sabanjure ndadekake
kompatibel; aturan pungkasan
njamin yen relasi iki ngemot
kontraksi-$\beta$.

Miturut intuisi, loro suku
ekuivalen-$\beta$ yen lan mung yen
siji bisa diowahi dadi sijine
kanthi nol utawa luwih langkah
kontraksi-$\beta$ utawa
``walikan'' kontraksi-$\beta$.
Walikan kontraksi-$\beta$
ditetepake supaya $M$
dikontraksi-$\beta$ kanthi
walikan dadi $N$ yen lan
mung yen $N$ dikontraksi-$\beta$
dadi $M$.

Kajaba aturan ing ndhuwur,
kita bakal nggedhekake relasi
kanthi aturan liyane lan
nglambangake relasi ekuivalensi
kang digedhekake iku minangka
$\equal[X]$, kanthi $X$
aturan tambahan kasebut.

\end{document}
